\documentclass[a4paper,12pt]{article}
\usepackage[T2A]{fontenc}
\usepackage[utf8]{inputenc}
\usepackage[russian]{babel}
\usepackage{amsmath,amssymb,mathtools}
\usepackage{geometry}
\geometry{top=23mm,bottom=25mm,left=24mm,right=24mm,headheight=14pt}
\usepackage{microtype}
\usepackage{booktabs,tabularx,array}
\usepackage{enumitem}
\usepackage{tikz}
\usepackage{fancyhdr,hyperref,xcolor}
\usepackage{needspace}
\usepackage{parskip}
\definecolor{accent}{HTML}{007C81}
\definecolor{err}{HTML}{A83F3F}
\definecolor{ink}{HTML}{303536}
\hypersetup{colorlinks=true,linkcolor=accent,urlcolor=accent,pdfauthor={Лёвин Артём Александрович},pdftitle={Ревью домашней работы — Марк Гукин}}
\setlength{\parindent}{0pt}
\setlength{\parskip}{0.55em}
\setlist[itemize]{leftmargin=1.6em,itemsep=0.25em,topsep=0.2em}
\pagestyle{fancy}
\fancyhf{}
\fancyhead[L]{\small\color{accent}Ревью домашней работы}
\fancyhead[R]{\small\color{ink}Марк Гукин}
\fancyfoot[C]{\scriptsize Лёвин Артём Александрович $\cdot$ эксклюзивно для Марка Гукина\quad|\quad\thepage}
\renewcommand{\headrulewidth}{0.25pt}
\renewcommand{\footrulewidth}{0pt}
\newcommand{\topic}[1]{{\color{accent}\bfseries #1}}
\newcommand{\mistake}[1]{{\color{err}\bfseries Ошибка:} #1}
\newcommand{\keyidea}[1]{{\color{accent}\bfseries Ключевая идея:} #1}
\newcommand{\summary}[1]{{\color{accent}\bfseries Итог:} #1}
\newcommand{\reviewsection}[1]{\vspace{0.55em}{\large\color{accent}\bfseries #1}\par\vspace{-0.3em}}
\begin{document}
\thispagestyle{empty}
\begin{center}
  \vspace*{38mm}
  {\color{accent}\large МАТЕМАТИКА}\par
  \vspace{15mm}
  {\LARGE Ревью домашней работы}\par
  \vspace{9mm}
  {\Large\color{accent}Комбинаторика и вероятность}\par
  \vspace{18mm}
  {\normalsize Проверка задач с подробным разбором ошибок}\par
  \vspace{29mm}
  {\large Марк Гукин}\par
  \vspace{13mm}
  {17 ноября 2026 г.}\par
  \vfill
  {\small Лёвин Артём Александрович}\par
  {\small эксклюзивно для Марка Гукина}\par
  \vspace{17mm}
  \begin{tikzpicture}
    \draw[accent,line width=0.75pt] (-4,0) .. controls (-2.85,0.9) and (-1.7,-0.4) .. (0,0)
      .. controls (1.7,0.4) and (2.75,-0.85) .. (4,0);
  \end{tikzpicture}
\end{center}
\clearpage
\phantomsection
\section*{Сводная таблица ошибок}
В таблицу включены задания, в которых обнаружены ошибки рассуждений или отсутствует окончательное решение. Верные задания в разборе опущены.
\par\medskip
\normalsize
\renewcommand{\arraystretch}{1.5}
\begin{tabularx}{\linewidth}{@{}p{10mm}p{25mm}X p{25mm}p{26mm}@{}}
\toprule
№ & Тема & Суть ошибки & Ваша запись & Правильный ответ\\
\midrule
\hyperref[task:9]{\textbf{9}} & Монеты & Неверно подсчитаны позиции орлов и записано ложное равенство дробей. & $5/16$ с вопросом & $5/16$\\
\hyperref[task:10]{\textbf{10}} & Числа и делимость & Неверное число всех исходов; подсчёт кратных четырём не завершён. & $96$; ответа нет & $6/25$\\
\bottomrule
\end{tabularx}
\normalsize
\par\vspace{13mm}
\textbf{Как работать с разбором.} Сначала найдите первый неверный переход в своей записи, затем восстановите правильное решение. После этого самостоятельно решите ровно одну задачу на закрепление в конце каждого разбора.
\clearpage
\thispagestyle{empty}
\vspace*{0.35\textheight}
\begin{center}
  {\color{accent}\LARGE Разбор ошибок}
\end{center}
\clearpage
\phantomsection\label{task:9}
\section*{Задание №\,9. Ровно три орла}
\topic{Тема:} независимые испытания и подсчёт сочетаний.

\reviewsection{Текст задания}
Честную монету независимо подбрасывают пять раз. Найдите вероятность того, что орёл выпадет ровно три раза.

\reviewsection{Ваше решение}
В записи читается указание на пять благоприятных исходов, затем дробь
\[
   \frac{5}{2^5}=\frac{5}{16}\,?
\]
Знак вопроса сохранён по смыслу записи: окончательная уверенность в ответе отсутствует.

\reviewsection{Ваш ответ}
\(\frac{5}{16}\) с вопросительным знаком. Численно это верный ответ, однако ход подсчёта содержит ошибки.

\reviewsection{Где именно возникает ошибка}
\mistake{Неверно определено количество способов получить ровно три орла.} Указание на пять возможных положений даёт слишком мало вариантов: орлы занимают \emph{три} из пяти мест, и выбрать нужно всю тройку позиций. Кроме того, переход от \(5/2^5\) к \(5/16\) арифметически неверен.

\textbf{Последний правильный шаг:} пять подбрасываний дают \(2^5\) равновозможных последовательностей, поскольку каждый раз есть два исхода.

\textbf{Первый неправильный шаг:} число благоприятных последовательностей принимается за \(5\).

\reviewsection{Почему это неверно}
Можно начать с пяти возможных положений первого орла, но после выбора одного положения остаётся выбрать ещё два. При прямом последовательном выборе одна и та же тройка пересчитывается несколько раз. Поэтому используется число \emph{сочетаний} трёх позиций из пяти. Также \(2^5=32\), поэтому \(5/2^5=5/32\), а дробь \(5/16\) так получить нельзя.

\keyidea{Сначала считаем все равновозможные последовательности, затем выбираем позиции ровно трёх орлов.}

\reviewsection{Исправление от последнего правильного шага}
Число всех исходов \(2^5=32\) сохраняем. Число благоприятных исходов пересчитываем через сочетания:
\[
  C_5^3=\frac{5\cdot4\cdot3}{3\cdot2\cdot1}=10.
\]
\needspace{9\baselineskip}
\reviewsection{Полное правильное решение}
\begin{enumerate}[leftmargin=1.75em,itemsep=.25em]
\item Каждое из пяти независимых подбрасываний даёт два исхода, поэтому всех последовательностей \(2^5=32\).
\item Чтобы орёл выпал ровно три раза, определяем три его позиции. Остальные две позиции автоматически займёт решка.
\item Выбрать три позиции из пяти можно \(C_5^3=10\) способами.
\item Все последовательности равновероятны, значит
\[
  P=\frac{10}{32}=\boxed{\frac{5}{16}}.
\]
\end{enumerate}

\reviewsection{Проверка}
Для симметричной монеты вероятность каждого конкретного расположения равна \((1/2)^5=1/32\). Таких расположений десять: \(10\cdot1/32=5/16\). Полученная вероятность меньше единицы и согласуется с прямым подсчётом.

\summary{Конечная дробь была записана верно, но две ошибки рассуждения необходимо устранить, чтобы результат был обоснован.}

\reviewsection{Задача на закрепление}
Честную монету независимо подбрасывают шесть раз. Найдите вероятность того, что орёл выпадет ровно четыре раза.
\clearpage
\phantomsection\label{task:10}
\section*{Задание №\,10. Шестизначные числа}
\topic{Тема:} перестановки с ограничением на первую цифру; признак делимости на 4.

\reviewsection{Текст задания}
Из цифр \(0,1,2,3,4,5\) случайно составляют шестизначное число, используя каждую цифру ровно один раз. Начальный нуль запрещён. Найдите вероятность того, что полученное число делится на 4.

\reviewsection{Ваше решение}
В записи видны попытки перемножить количества вариантов по разрядам и пометка «общ. = 96». Окончательная дробь вероятности отсутствует. Часть зачёркнутых промежуточных пометок недостаточно надёжна для подробной реконструкции, поэтому их смысл не предполагаем.

\reviewsection{Ваш ответ}
Не указан.

\reviewsection{Где именно возникает ошибка}
\mistake{Для шестизначных чисел из шести различных цифр общее количество \(96\) неверно.} В первом разряде можно поставить любую из пяти ненулевых цифр; во всех остальных разрядах оставшиеся пять цифр переставляются произвольно.

\textbf{Последний достоверно правильный замысел:} считать варианты выбора цифр по разрядам с помощью произведения.

\textbf{Первый достоверно неверный результат:} пометка «общ. = 96». Верное значение общего количества --- \(600\).

\reviewsection{Почему это неверно}
Количество вариантов при заполнении очередного разряда зависит от того, какие цифры уже использованы. После выбора первой ненулевой цифры остаются ровно пять различных цифр, включая нуль, который на этом этапе разрешён. Значит, число вариантов заполнения пяти последних разрядов равно \(5!=120\). При подсчёте чисел, кратных четырём, дополнительно важно проверять \emph{две последние цифры}: именно они определяют делимость всего числа на 4.

\keyidea{Разделяем подсчёт всех допустимых чисел и подсчёт чисел с разрешёнными окончаниями. Для разных окончаний количество первых четырёх цифр может отличаться.}

\reviewsection{Исправление от правильного замысла}
Для общего числа исходов получаем
\[
 N=5\cdot5!=5\cdot120=600.
\]
Теперь определим благоприятные окончания. Двузначная запись из цифр \(0,1,2,3,4,5\) с разными цифрами делится на 4 в следующих случаях:
\[
 04,\quad12,\quad20,\quad24,\quad32,\quad40,\quad52.
\]
Начальный нуль в \emph{двузначном окончании} допустим: например, число, оканчивающееся на \(04\), может быть шестизначным.

\reviewsection{Полное правильное решение}
Разбиваем семь окончаний на две группы.

\textbf{Случай 1: нуль находится в окончании.}
Это окончания \(04,20,40\) --- три варианта. Впереди остаются четыре ненулевые цифры. Их можно переставить \(4!=24\) способами. Получаем
\[
   N_1=3\cdot24=72.
\]
\textbf{Случай 2: нуля в окончании нет.}
Это окончания \(12,24,32,52\) --- четыре варианта. Среди оставшихся четырёх цифр есть нуль, поэтому для первой позиции есть ровно три ненулевых варианта, а для трёх остальных --- \(3!=6\) перестановок. Значит,
\[
   N_2=4\cdot3\cdot3!=4\cdot18=72.
\]
Всего чисел, кратных четырём,
\[
 N_{\text{благ}}=72+72=144.
\]
Каждое из \(600\) допустимых чисел равновероятно, поэтому
\[
 P=\frac{144}{600}=\boxed{\frac{6}{25}}.
\]

\reviewsection{Проверка}
Подсчёт двух случаев независим и полон: каждое допустимое окончание содержится ровно в одной группе. Благоприятных чисел \(144\), что меньше общего количества \(600\); дробь сокращается делением числителя и знаменателя на \(24\).

\summary{Правильная вероятность равна \(6/25\). Критические элементы решения: первая цифра не равна нулю, признак делимости на 4 и отдельный учёт окончаний, в которых встречается нуль.}

\reviewsection{Задача на закрепление}
Из цифр \(0,1,2,3,4\) наугад составляют пятизначное число, используя каждую цифру ровно один раз. Начальный нуль запрещён. Найдите вероятность того, что полученное число делится на 4.
\end{document}
